\documentclass[11pt,a4paper]{article}
\usepackage[utf8]{inputenc}
\usepackage{graphicx}
\usepackage[left=2.5cm,top=3cm,right=2.5cm,bottom=3cm,bindingoffset=0.5cm]{geometry}
\usepackage{AEDLogica, AEDEspecificacion, AEDTADs}
\usepackage{caratula}
\usepackage{hyperref}
\newlength{\miSangria}
\setlength{\miSangria}{2em}

\titulo{Trabajo práctico 1: Especificacion de Tads}
\subtitulo{Tad AirMiles}

\fecha{\today}

\materia{Algoritmos y Estructuras de Datos}
\grupo{Grupo: Arraydecuatroposiciones} 

\integrante{Javier M. Ramos}{288/20}{ramosjavier276@gmail.com}
\integrante{Nancy Ramos}{1206/24}{Ramosnancydaniela@outlook.com}
\integrante{Elian Carreño}{1420/23}{elianuba2004@gmail.com }
\integrante{Leandro Amodey}{618/24}{leandroamodey@gmail.com}
% \integrante{Apellido, Nombre2}{002/01}{email2@dominio.com} %


% Declaramos donde van a estar las figuras
% No es obligatorio, pero suele ser comodo
\graphicspath{{../static/}} 

% Asi pueden escribir nuevos comandos. 
% Este por ejemplo asegura q los nombres 
% que figuren con una tipografia diferenciada  
\newcommand{\Tipo}[1]{\mathsf{#1}}
% la sintaxis es \newcommand{\nombreDeLaMacro}[cantidadDeParametros]{Lo que va ser remplazado por el macro} 
\newcommand{\norm}[1]{\vert#1\vert}



\begin{document}
\maketitle
\section{Introducción}

El presente trabajo practico es una especificacion formal para un sistema de recompensas basado en millas, como ocurre en muchas aerolineas reales.
\\
** VER CUANTO DETALLE DE NUESTRA INTERPRETACION PONEMOS EN LA INTRO **

\section{Renombres de Tipos y Estructuras}

\noindent enum Categoría \textbf{ES} \{Bronce, Plata,Oro, Platino  \}

\vspace{0.5cm}

\noindent Transacciones \textbf{ES} Struct $\tupla{
		UsuarioOrigen: \Tipo{\Z}, \\
		UsuarioDestino: \Tipo{\Z},
		TipoDeTransaccion: \Tipo{\Z},
		Millas: \Tipo{\R}, \\
		Distancia: \Tipo{\Z},
		NombrePromo: \Tipo{\str},
		FactorPromo: \Tipo{\R}
	}$

\noindent \comentario{Tipo de transacciones es algo que yo pongo a mano creo, ejemplo 0:Transaccion por referidor, 1: Transferencia, 2: Canjeo de puntos}

\vspace{2em}

\section{Especificación del TAD AirMiles}
\begin{tad}{Airmiles}


\comentario{Decidimos dejar en usuarios la categoria por mas que se pueda calcular por fuera. Las promociones tienen a un diccionario que adentro tiene un struct }
\\
\noindent \comentario{La convencion utilizada para transacciones es:
	\\
	1) Si se hace un vuelo, el TipoDeTransaccion es 0 y el usuario de destino es -1
	\\
	2) Si se hace una transferencia el TipoDeTransaccion es 1. La distancia en este caso se pasa como 0.
	\\
	3) Si se hace un canjeo el TipoDeTransaccion es 2 y el usuario de destino es -1, se modifica unicamente el usuario de origen, la distancia en este caso se pasa como 0.
}
\\
\comentario{En caso de las promociones, si la promocion no tiene un tope, por convencion tomamos millasRestantes = -1.}
\\
\obs{Usuarios}{\dict{id:\Tipo{\Z}}{CategoriaActual : Categoria}}
\\
\noindent \obs{Historial}{\seq{Transacciones}}

\texttt{obs} Promociones : dict \{ \\
\hspace*{1.5em} NombrePromocion: $\str$, \\
\hspace*{1.5em} $\struct{\Tipo{FactorDePromocion} : \R,\Tipo{topeDePromocion}: \R, \Tipo{millasRestantes}: \R}$ \\
\}


\phantomsection\hypertarget{sec:crearSistema}{}\label{sec:crearSistema}
\begin{proc}{crearSistema}
	{}
	{AirMiles}
	\requiere{\Tipo{True}}
	\aseguraLargo{ \norm{res.Usuarios} = 0 \land
		\\
		\norm{res.Historial} = 0 \land
		\\
		res.promociones = setKey(\dict{}{},"NoPromo",\struct{1,0,-1}) }
\end{proc}

\phantomsection\hypertarget{sec:registrarUsuario}{}\label{sec:registrarUsuario}
\begin{proc}{registrarUsuario}
	{
		\Inout AM: Airmiles,
		\In id: \Tipo{\Z}
	}
	{}
	\comentario{Actualiza el diccionario de usuarios con  nuevo usuario} \\
	\requiere{AM = AM0 \land id \notin claves(AM.Usuarios) \land id > 0}
	\aseguraLargo{ AM.Usuarios = SetKey(AM0.Usuarios, id, Bronce) \land \\
		AM.Historial = AM0.Historial \\
		\land AM.Promociones = AM0.Promociones
	}
\end{proc}


\phantomsection\hypertarget{sec:registrarUsuarioConReferidor}{}\label{sec:registrarUsuarioConReferidor}
\begin{proc}{registrarUsuarioConReferidor}
	{
		\Inout AM: Airmiles,\\
		\In idReferidor: \Tipo{\Z},
		\In idUsuario: \Tipo{\Z}
	}
	{}
	\requiere{AM0 = AM \land idReferidor \in claves(AM.Usuarios) \land idUsuario \notin claves(AM.Usuarios) \land idUsuario > 0}
	\aseguraLargo{
		% Creo el usuario
		AM.Usuarios = SetKey(AM0.Usuarios,idUsuario, Bronce) \land
		\\
		\hyperref[sec:EsNuevaTransaccion]{EsNuevaTransaccion}(AM,AM0,idReferidor, idUsuario, 0, \\
		500,0,"NoPromo",1)\land
		\\
		AM.Promociones = AM0.Promociones
	}
\end{proc}
% ===============Aca comienza la funcion de acumular millas===============

\noindent \comentarioLargo{Acumula millas al usuario si la categoria del mismo no es platino, en caso de que el nombre de la promo no alcance para completar el viaje del usuario, se crean dos transacciones, una con las millas restantes de la promo y otra con "NoPromo" para completar el viaje. }
\phantomsection\hypertarget{sec:AcumularMillas}{}\label{sec:AcumularMillas}
\begin{proc}{AcumularMillas}
	{
		\Inout AM: \Tipo{Airmiles},
		\In idUsuario: \Tipo{\Z},
		\In distancia: \Tipo{\R},
		\In nombrePromo: \Tipo{\str},
	}
	{}

	\requiere{AM = AM0 \land idUsuario \in claves(AM.Usuarios) \land nombrePromo \in claves(AM.Promociones) \land distancia > 0 }
	\aseguraLargo{
		\comentario{Si el tope de las millas de la promo supera la distancia se hace una sola transaccion}
		\\
		AM0.Promociones\lbrack nombrePromo \rbrack.millasRestantes \geq
		\\
		\hyperref[sec:MillasCalculadas]{MillasCalculadas}(distancia, \hyperref[sec:DarFactorCategoria]{DarFactorCategoria}(AM0.Historial,idUsuario),
		\\ AM0.Promociones[nombrePromo].factorDePromocion) 
		\lor 
		\\
		\\
		\comentario{Si la promo no tiene tope lo manejamos a parte}
		\\
		AM0.Promociones\lbrack nombrePromo \rbrack.millasRestantes = -1
		\implicaLuego
		\\
		\comentario{Se crea la transaccion}
		\\
		\hyperref[sec:EsNuevaTransaccion]{EsNuevaTransaccion}(AM,AM0,idUsuario,-1,0,
		\\
		\hyperref[sec:MillasCalculadas]{MillasCalculadas}(distancia, \hyperref[sec:DarFactorCategoria]{DarFactorCategoria}(AM0.Historial,idUsuario),
		\\ AM0.Promociones[nombrePromo].factorDePromocion) ,distancia, \\
		nombrePromo, AM0.Promociones[nombrePromo].FactorDePromocion) \land\\
		\\
		\comentario{Se restan las millas de la promocion unicamente para promos que no tiene tope}
		\\
		AM0.Promociones\lbrack nombrePromo \rbrack.millasRestantes > 0 \implicaLuego
		\\ 
		AM.Promociones = setKey(AM0.Promociones, nombrePromo,\\
		AM0.Promociones[nombrePromo].millasRestantes - \hyperref[sec:MillasCalculadas]{MillasCalculadas}\\(distancia, \hyperref[sec:DarFactorCategoria]{DarFactorCategoria}(AM0.Historial, idUsuario), \\
		AM0.Promociones[nombrePromo].factorDePromocion)
		\lor
		\\
		\\
		AM0.Promociones\lbrack nombrePromo \rbrack.millasRestantes <
		\\
		\hyperref[sec:MillasCalculadas]{MillasCalculadas}(distancia, \hyperref[sec:DarFactorCategoria]{DarFactorCategoria}(AM0.Historial, idUsuario),\\
		AM0.Promociones[nombrePromo].factorDePromocion) \implica
		\\
		\\
		\comentario{Se hace 1 transaccion con las millas restantes de promo}
		\\
		\hyperref[sec:EsNuevaTransaccion]{EsNuevaTransaccion}(AM,AM0,idUsuario,-1,1,
		\\
		AM0.Promociones[nombrePromo].millasRestantes, distancia,
		\\ nombrePromo, AM0.Promociones[nombrePromo].FactorDePromocion) \yLuego
		\\
		\\
		\comentario{Se hace otra transaccion con noPromo para completar}
		\\
		\hyperref[sec:EsNuevaTransaccion]{EsNuevaTransaccion}(AM,AM0,idUsuario,-1,1,
		\\
		MillasACompletar(AM0.Promociones,"NoPromo",distancia,
		\\\hyperref[sec:DarFactorCategoria]{DarFactorCategoria}(AM0.Historial,idUsuario), 1)
		\\distancia, "NoPromo", 1) \yLuego
		% Luego de hacer la transaccion tengo que verificar si se modifico la categoría del usuario
		\\
		\comentario{Se puede indefinir si no le pasamos un usuario que este en el sistema}
		\\
		\hyperref[sec:SubioCategoria]{SubioCategoria}(idUsuario,AM.Historial,AM0.Historial) \implicaLuego
		\\
		AM.Usuarios = SetKey(AM0.Usuarios,idUsuario,
		\\
		\hyperref[sec:DarCategoria]{DarCategoria}(AM0.Historial, idUsuario)) \land
		\\
		\comentario{Despues de completar el viaje se verifica si la promocion termino}
		\\
		AM0.Promociones[nombrePromo].millasRestantes = 0 \implicaLuego
		\\
		\hyperref[sec:SeConsumioLaPromocion]{SeConsumioLaPromocion}(nombrePromo, AM, AM0)
	}
\end{proc}

\vspace{1cm}
\phantomsection\hypertarget{sec:CanjearMillas}{}\label{sec:CanjearMillas}
\begin{proc}{CanjearMillas}
	{\Inout AM: AirMiles, \In MillasCanjeadas : \Tipo{\R}, \In IdUsuario : \Tipo{\Z}}
	{}
	\requiereLargo{ AM=AM0 \yLuego MillasCanjeadas > 0 \land IdUsuario \in claves(AM.Usuarios)}
	\aseguraLargo{
		\IfThenElse{ \hyperref[sec:MillasDisponibles]{MillasDisponibles}(IdUsuario, AM0.Historial) \geq MillasCanjeadas}
		{\\EsNuevaTransaccion(AM,AM0,IdUsuario,-1,2,MillasCanjeadas,0,"NoPromo",1)}
		{\\AM.Historial =  AM0.Historial} \\
		\land AM.Usuarios = AM0.Usuarios \\
		\land AM.Promociones = AM0.Promociones
	}
\end{proc}

\vspace{1cm}
\phantomsection\hypertarget{sec:TransferirMillas}{}\label{sec:TransferirMillas}
\begin{proc}{TransferirMillas}
	{
		\Inout AM: Airmiles,\\
		\In idUsuarioPagador: \Tipo{\Z},
		\In idUsuarioBeneficiario: \Tipo{\Z},
		\In CantMillasTrans: \Tipo{\R}
	}
	{}
	\requiere{AM0 = AM \land idUsuarioPagador \in claves(AM.usuarios)\\
		\land idUsuarioBeneficiario \in claves(AM.usuarios)\\
		\land idUsuarioPagador \neq idUsuarioBeneficiario\\
		\land \hyperref[sec:MillasDisponibles]{MillasDisponibles}(AM, idUsuarioPagador) \geq CantMillasTrans
	}
	\aseguraLargo{\hyperref[sec:EsNuevaTransaccion]{EsNuevaTransaccion}( AM, AM0, idUsuarioPagador, idUsuarioBeneficiario
	\\,1, CantMillasTrans,0, NoPromo, 1) \yLuego 
		\\ \hyperref[sec:SubioCategoria]{SubioCategoria}(idUsuarioBeneficiario) \land 
		\\
		AM.Promociones = AM0.Promociones \land
		\\
		AM.Usuarios = AM0.Usuarios
	}
	\noindent \comentarioLargo{Revisado}

\end{proc}

\vspace{2cm}
\phantomsection\hypertarget{sec:ConsultarSaldo}{}\label{sec:ConsultarSaldo}
\begin{proc}{ConsultarSaldo}
	{\In AM: AirMiles, \In IdUsuario : \Tipo{\Z}}
	{\Tipo{\Z}}
	\requiereLargo{IdUsuario \in claves(AM.Usuarios) }

	\aseguraLargo{
		res = \hyperref[sec:MillasDisponibles]{MillasDisponibles}(IdUsuario,AM.Historial)
	}
\end{proc}
\comentario{El auxiliar MillasDisponibles es el que esta en CanjearMillas. Ver despues si reusamos ese para todo u algun otro, pero hay que unificar.}

\vspace{1cm}
\phantomsection\hypertarget{sec:ConsultarCategoria}{}\label{sec:ConsultarCategoria}
\begin{proc}{ConsultarCategoria}
	{
		\In AM: Airmiles,
		\In idUsuario: \Tipo{\Z},
	}
	{\Tipo{categoria}}
	\requiere{idUsuario \in claves(AM.usuarios)}
	\asegura{res = AM.Usuarios[idUsuario]}

\end{proc}

\vspace{1cm}
\phantomsection\hypertarget{sec:RankingMillas}{}\label{sec:RankingMillas}
\begin{proc}{RankingMillas}
	{\In AM: Airmiles, \In n : \Tipo{\Z}}
	{\Tipo{\seq{\Tipo{\Z}}}}
	\requiereLargo{ n > 0}

	\aseguraLargo{ (|AM.Usuarios| > n  \land_L |res| = n) \\ \lor (|AM.Usuarios| \leq n  \land_L |res| = |AM.Usuarios|)
		\\ \land_L \hyperref[sec:perteneceUsuarios]{perteneceUsuarios}(AM.Usuarios, res)
		\\ \land_L \hyperref[sec:sinRepetidos]{sinRepetidos}(res)
		\\ \land_L \hyperref[sec:ordenDecreciente]{ordenDecreciente}(AM.Historial, res)
		\\ \land_L \hyperref[sec:losMejoresCandidatos]{losMejoresCandidatos}(AM.Usuarios, AM.
		Historial, res)
	}
\end{proc}

\vspace{0.5cm}
\phantomsection\hypertarget{sec:ResumenTransacciones}{}\label{sec:ResumenTransacciones}
\comentario{Falta implementar el auxiliar MillasTransferidasDelUsuario.}
\\
\begin{proc}{ResumenTransacciones}
	{\In AM: Airmiles, \In id : \Tipo{\Z}}
	{\tupla{\Tipo{\R}, \Tipo{\R}, \Tipo{\R}}}
	\requiereLargo{ id \in claves(AM.Usuarios)}

	\aseguraLargo{ res_{0} = \hyperref[sec:MillasObtenidasPorVuelo]{MillasObtenidasPorVuelo}(AM.Historial, id) \yLuego
		\\
		res_{1} = \hyperref[sec:MillasDescontadasPorCanjeo]{MillasDescontadasPorCanjeo}(AM.Historial, id) \yLuego
		\\
		res_{2} = \hyperref[sec:MillasTransferidasYRecibidas]{MillasTransferidasYRecibidas}(AM.Historial, id)
	}
\end{proc}

\phantomsection\hypertarget{sec:MillasTransferidasYRecibidas}{}\label{sec:MillasTransferidasYRecibidas}
\auxLargo{MillasTransferidasYRecibidas}{Historial:\Tipo{\seq{Transaccion}}, id: \Tipo{\Z}}{\Tipo{\R}}
{
	\sum\limits_{0}^{\norm{Historial}-1}\IfThenElse{EsTransferenciaEnviada}{-Historial[i].Millas}
	\\{\IfThenElse{EsTransferenciaRecibida}{Historial[i].Millas}{0}}
}
\predLargo{EsTransferenciaRecibida}
{Transaccion: \Tipo{Transaccion}, idDestino: \Tipo{\Z}, Usuarios: \Tipo{\dict{\Z}{Categoria}}}
{
	Transaccion.UsuarioOrigen \in claves(Usuarios) \land
	\\
	Transaccion.TipoDeTransaccion = 1 \land
	\\
	Transaccion.UsuarioDestino = idDestino
}

\predLargo{EsTransferenciaEnviada}
	{Transaccion: \Tipo{Transacciones}, idOrigen: \Tipo{\Z}, Usuarios: \Tipo{\dict{\Z}{Categoria}}}
	{			
		Transaccion.UsuarioOrigen = idOrigen \land
		\\
		Transaccion.TipoDeTransaccion = 1 \land
		\\
		Transaccion.UsuarioDestino \in claves(Usuarios) 
	}

\vspace{1cm}
\phantomsection\hypertarget{sec:IniciarPromocion}{}\label{sec:IniciarPromocion}
\begin{proc}{IniciarPromocion}
	{
		\Inout AM: \Tipo{Airmiles},
		\In nombrePromo: \Tipo{\seq{char}},
		\In factorPromo: \Tipo{\R},
	}
	{}
	\comentario{Discutir el tercer elemento de la lista [factorPromo,0,0].
		NoPromo esta siempre en el sistema, no se puede pisar. Discutir si seria != noPromo o si directamente no tiene que pertenecer (podria existir y pisarse el valor?)}
	\\
	\requiereLargo{
		AM = AM0 \land
		factorPromo \geq 1 \land
		nombrePromo \notin claves(AM0.Promociones)}
	\aseguraLargo{
		AM.Usuarios = AM0.Usuarios \land
		\\
		AM.Historial = AM0.Historial \land
		\\
		AM.Promociones = SetKey(AM0.Promociones, nombrePromo, \struct{factorPromo,0,-1})
	}
\end{proc}
\vspace{5.0cm}

\phantomsection\hypertarget{sec:IniciarPromocionConTope}{}\label{sec:IniciarPromocionConTope}
\begin{proc}{IniciarPromocionConTope}
{
	\Inout AM: \Tipo{Airmiles},
	\In nombrePromo: \Tipo{\str},
	\In factorPromo: \Tipo{\R},
	\In TopePromo: \Tipo{\R}
}
{}
\requiereLargo{
	AM = AM0 \land factorPromo \geq 1 \land nombrePromo \notin
	claves(AM.Promociones) \land TopePromo > 0
}
\aseguraLargo{
	AM.Usuarios = AM0.Usuarios \land
	\\
	AM.Historial = AM0.Historial \land
	\\
	AM.Promociones = SetKey(AM0.Promociones, nombrePromo,
	\\
	\struct{factorPromo,topePromo,topePromo})
}
\end{proc}
% Asegura{ AM.Usuarios = AM0.Usuarios  ∧ 
% 		   AM.Historial = AM0.Historial ∧
% 		   AM.Promociones = SetKey(AM0.Promociones, nombrePromo, [factorPromo,0,0] 		
% 		 }

\vspace{1.0cm}

\phantomsection\hypertarget{sec:FinalizarPromocion}{}\label{sec:FinalizarPromocion}
\begin{proc}{FinalizarPromocion}
	{
		\Inout AM: \Tipo{Airmiles},
		\In nombrePromo: \Tipo{\seq{char}}
	}
	{}
	\requiereLargo{
		AM = AM0 \land nombrePromo \neq "NoPromo" \yLuego
		\\
		nombrePromo \in claves(AM.Promociones)
	}
	\aseguraLargo{
		AM.Usuarios = AM0.Usuarios \land
		\\
		AM.Historial = AM0.Historial \land
		\\
		AM.Promociones = DelKey(AM0.Promociones, nombrePromo)
	}
\end{proc}

\vspace{1.0cm}

\phantomsection\hypertarget{sec:ElQueTieneMasMillas}{}\label{sec:ElQueTieneMasMillas}
\begin{proc}{ElQueTieneMasMillas}
	{
		\In AM: \Tipo{Airmiles}
	}
	{\Tipo{\Z}}
	\requiere{ \norm{AM.Usuarios} > 0}
	\aseguraLargo{
		\existe{i}{\Tipo{\Z}}{i \in claves(AM.Usuarios) \yLuego res = i } \land
		\\
		\paraTodo{j}{\Tipo{\Z}}{j \in claves(AM.Usuarios) \implicaLuego \hyperref[sec:MillasTotales]{MillasTotales}(res,AM.Usuarios) \geq \hyperref[sec:MillasTotales]{MillasTotales}(j,AM.Usuarios)}
	}
\end{proc}

\vspace{1.0cm}



\phantomsection\hypertarget{sec:ElMasCanjeador}{}\label{sec:ElMasCanjeador}
\begin{proc}{ElMasCanjeador}
	{
		\In AM: \Tipo{Airmiles}
	}
	{\Tipo{\Z}}
	\requiere{
		\norm{AM.Usuarios} > 0 \land \norm{AM.Historial} > 0
	}
	\aseguraLargo{
		\existe{id}{\Tipo{\Z}}{id \in claves(AM.Usuarios) \yLuego res = id
			\\
			\land \paraTodo{id2}{\Tipo{\Z}}{id2 \in claves(AM.Usuarios) \implicaLuego
				\\
				\hyperref[sec:MillasPerdidas]{MillasPerdidas}(AM.Historial,res,AM.Usuarios) \ge
				\\
				\hyperref[sec:MillasPerdidas]{MillasPerdidas}(AM.Historial,id2,AM.Usuarios)}}
	}
\end{proc}

\vspace{1.cm}
\comentario{Falta Implementar: ParesDeUsuariosValidos y EstanTodosLosParesExclusivos.}
\phantomsection\hypertarget{sec:paresTransferidoresExclusivos}{}\label{sec:paresTransferidoresExclusivos}
\begin{proc}{paresTransferidoresExclusivos}
		{\In AM: Airmiles}
		{\seq{\tupla{\Tipo{\Z}, \Tipo{\Z}}}}
		\requiere{True}

		\aseguraLargo{
            ParesDeUsuariosValidos(AM.Usuarios, res) \land_L
            \\
            \hyperref[sec:sinRepetidos]{sinRepetidos}(res) \land_L
            \\
            \hyperref[sec:NoHayPermutacion]{NoHayPermutacion}(res) \land_L
            \\
            \hyperref[sec:TieneReciprocoHistorial]{TieneReciprocoHistorial}(AM.Historial, res) \land_L
            \\
            \hyperref[sec:SonParejaExclusiva]{SonParejaExclusiva}(AM.Historial, res) \land_L
            \\
            EstanTodosLosParesExclusivos(
                AM.Usuarios,
                AM.Historial,
                res
            )
        }
	\end{proc}


\phantomsection\hypertarget{sec:perteneceUsuarios}{}\label{sec:perteneceUsuarios}
\predLargo{perteneceUsuarios}{usuarios:\Diccionario{\Tipo{\Z}}{categoria}, res: \seq{\Z}} {
	\paraTodo{i}{\Tipo{\Z}}{0 \leq i < \norm{res} \implica_L res[i] \in claves(usuarios)}
}

\phantomsection\hypertarget{sec:sinRepetidos}{}\label{sec:sinRepetidos}
\predLargo{sinRepetidos}{res: \seq{\Z}} {
	\paraTodo{x}{\Tipo{\Z}}{x \in res\implica_L \hyperref[sec:cantAparicion]{cantAparicion}(res,x) = 1}
}

\phantomsection\hypertarget{sec:cantAparicion}{}\label{sec:cantAparicion}
\auxLargo{cantAparicion}
{res: \Tipo{\seq{\Z}}, x: \Tipo{\Z}}
{\Tipo{\R}}
{
	\sum\limits_{i=0}^{\norm{res} - 1}
	\IfThenElse{
		res[i] = x
	}{
		1
	}{
		0
	}
}

\phantomsection\hypertarget{sec:ordenDecreciente}{}\label{sec:ordenDecreciente}
\predLargo{ordenDecreciente}{Historial:\seq{Transacciones}, res: \seq{\Z}} {
	\paraTodo{i}{\Tipo{\Z}}{0 \leq i < \norm{res}-1 \implica_L \\ millasDisponibles(res[i], Historial) >= millasDisponibles(res[i + 1], Historial)}
}
\\
\phantomsection\hypertarget{sec:losMejoresCandidatos}{}\label{sec:losMejoresCandidatos}
\predLargo{losMejoresCandidatos}{usuarios:\Diccionario{\Tipo{\Z}}{categoria}, \\ Historial:\seq{Transacciones}, res: \seq{\Z}} {
	\paraTodo{i}{\Tipo{\Z}}{0 \leq i < \norm{res} \land  res[i] \in claves(usuarios)  \implica_L  \\ \hyperref[sec:elMejorCandidato]{elMejorCandidato}(Usuarios, Historial, res, i) }
}
\\
\phantomsection\hypertarget{sec:elMejorCandidato}{}\label{sec:elMejorCandidato}
\predLargo{elMejorCandidato}{usuarios:\Diccionario{\Tipo{\Z}}{categoria}, \\ Historial:\seq{Transacciones}, res: \seq{\Z}, i: \Tipo{\Z}} {
	\paraTodo{x}{\Tipo{\Z}}{ (x \in claves(usuarios) \land \\ millasDisponibles(x,Historial) \leq millasDisponibles(res[i],Historial))\implica_L  \\ (\existe {j}{\Tipo{\Z}}{0 \leq j < i } \land_L res[j] = x) }
}

\vspace{1cm}
\phantomsection\hypertarget{sec:SubioCategoria}{}\label{sec:SubioCategoria}
\predLargo
{SubioCategoria}
{idUsuario: \Tipo{\Z},Historial: \Tipo{Transaccion},Historial0: \Tipo{Transaccion}}
{\hyperref[sec:DarCategoria]{DarCategoria}(Historial0,idUsuario) \neq \hyperref[sec:DarCategoria]{DarCategoria}(Historial,idUsuario)}


\vspace{1cm}
\phantomsection\hypertarget{sec:MillasCalculadas}{}\label{sec:MillasCalculadas}
\auxLargo
{MillasCalculadas}
{distancia:\Tipo{\Z}, Factorcategoria: \Tipo{\R}, factorPromo: \Tipo{\R}}
{\R}
{distancia * Factorcategoria * factorPromo}

\vspace{1cm}
\noindent \comentarioLargo{Con respecto a las transacciones, deberiamos limitar esto no? cuando el usuario es platino creo que no tendria sentido crearle transferencias para que gane millas, se supone que no podria. Por el contrario si puede crear transacciones que de millas}
\phantomsection\hypertarget{sec:EsNuevaTransaccion}{}\label{sec:EsNuevaTransaccion}
\predLargo{EsNuevaTransaccion}{
	AM,AM0 : \Tipo{AirMiles},idUsuario: \Tipo{\Z}, idReferidor: \Tipo{\Z}, \\TipoDeTransaccion: \Tipo{\Z}, Millas: \Tipo{\R}, distancia: \Tipo{\Z}, NombrePromo: \Tipo{\str}, FactorPromo: \Tipo{\R}
}
{AM.Historial = AM0.Historial ++ \lbrack \tupla{idUsuario, idReferidor, \\ TipoDeTransaccion, Millas,NombrePromo, FactorPromo} \rbrack}
{}

\vspace{1cm}
\phantomsection\hypertarget{sec:DarFactorCategoria}{}\label{sec:DarFactorCategoria}
\auxLargo{DarFactorCategoria}
{Historial: \Tipo{\seq{Transacciones}}, idUsuario: \Tipo{\Z}}
{\Tipo{\R}}
{
	\IfThenElse{
		\hyperref[sec:MillasTotales]{MillasTotales}(Historial, idUsuario) < 10000
	}{
		1.0
		\\
	}{
		\IfThenElse{
			\hyperref[sec:MillasTotales]{MillasTotales}(Historial, idUsuario) < 50000
		}{
			1.5
			\\
		}{
			\IfThenElse{
				\hyperref[sec:MillasTotales]{MillasTotales}(Historial, idUsuario) < 1000000
			}{
				2.0
			}{
				0.0
			}
		}
	}
}

\vspace{0.5cm}
\phantomsection\hypertarget{sec:DarCategoria}{}\label{sec:DarCategoria}
\auxLargo{DarCategoria}
{Historial: \Tipo{\seq{Transacciones}}, idUsuario: \Tipo{\Z}}
{\Tipo{Categoria}}
{
	\IfThenElse{
		\hyperref[sec:MillasTotales]{MillasTotales}(Historial, idUsuario) < 10000
	}{
		Bronce
		\\
	}{
		\IfThenElse{
			\hyperref[sec:MillasTotales]{MillasTotales}(Historial, idUsuario) < 50000
		}{
			Plata
			\\
		}{
			\IfThenElse{
				\hyperref[sec:MillasTotales]{MillasTotales}(Historial, idUsuario) < 1000000
			}{
				Oro
			}{
				Platino
			}
		}
	}
}


\vspace{0.5cm}
\noindent \comentario{En caso de que el usuario haga parte del viaje con promo y otra parte sin promo esta funcion busca esa diferencia}
\\
\phantomsection\hypertarget{sec:MillasACompletar}{}\label{sec:MillasACompletar}
\auxLargo{MillasACompletar}
{Promociones: \Tipo{\Diccionario{\str}{struct}},
	\\nombrePromo: \Tipo{\str}, distancia: \Tipo{\Z}, FactorCategoria: \Tipo{\R}}
{\R}
{
	\hyperref[sec:MillasCalculadas]{MillasCalculadas}(distancia,factorCategoria,\\
	Promociones[nombrePromo].factorDePromocion)-
	\\
	Promociones[nombrePromo].millasRestantes
}


\vspace{0.5cm}
\comentario{¿puedo chequear como va la sumatoria mientras va recorriendo?}\\
\phantomsection\hypertarget{sec:MillasTotales}{}\label{sec:MillasTotales}
\auxLargo{MillasTotales}
{IdUsuario : \Tipo{\Z}, Historial : \seq{Transacciones}}
{\Tipo{\Z}}
{res = \sum\limits_{i=0}^{\lvert Historial \rvert - 1} \IfThenElse{\hyperref[sec:EsTransaccionDelUsuario]{EsTransaccionDelUsuario}(Historial[i], IdUsuario \land \hyperref[sec:EsTransaccionSumadora]{EsTransaccionSumadora}(Historial[i], IdUsuario)) \\
		\land 1000000 > res}{Historial[i].Millas}{0}}



\vspace{0.5cm}
\comentarioLargo{Auxiliar para obtener las millas disponibles de un usuario. Recorre el historial de transacciones y de acuerdo al TipoDeTransaccion, al IdUsuario y las Millas del registro, suma o resta con un acumulador.}
\phantomsection\hypertarget{sec:MillasDisponibles}{}\label{sec:MillasDisponibles}
\auxLargo{MillasDisponibles}
{IdUsuario : \Tipo{\Z}, Historial : \seq{Transacciones}}
{\Tipo{\Z}}
{res = \sum\limits_{i=0}^{\lvert Historial \rvert - 1} \IfThenElse{\hyperref[sec:EsTransaccionDelUsuario]{EsTransaccionDelUsuario}(Historial[i], IdUsuario)}{\\\IfThenElse{\hyperref[sec:EsTransaccionSumadora]{EsTransaccionSumadora}(Historial[i], IdUsuario)\\}{Historial[i].Millas}{-Historial[i].Millas}}{0}}


\vspace{1.0cm}
\phantomsection\hypertarget{sec:EsTransaccionSumadora}{}\label{sec:EsTransaccionSumadora}
\predLargo{EsTransaccionSumadora}{Trans: Transacciones,idUsuario : \Tipo{\Z}}
{ (Trans.TipoDeTransaccion = 0) \lor \\(IdUsuario = Trans.UsuarioDestino \land Trans.TipoDeTransaccion = 1)}


\vspace{1.0cm}
\phantomsection\hypertarget{sec:EsTransaccionDelUsuario}{}\label{sec:EsTransaccionDelUsuario}
\pred{EsTransaccionDelUsuario}{Trans: Transacciones,idUsuario : \Tipo{\Z}}{\\Trans.UsuarioOrigen = idUsuario \lor Trans.UsuarioDestino = idUsuario}

\comentario{Valido que el IdUsuario este en alguno de los 2 campos. Si no esta en ninguno, entonces no debo sumar ni restar nada.}

\vspace{0.5cm}
\phantomsection\hypertarget{sec:SeConsumioLaPromocion}{}\label{sec:SeConsumioLaPromocion}
\predLargo
{SeConsumioLaPromocion}
{
	nombrePromo: \Tipo{\str},
	AM: Airmiles
	AM0: Airmiles
}
{AM.Promociones = DelKey(AM0.Promociones,nombrePromo)}

\vspace{0.5cm}

\phantomsection\hypertarget{sec:MillasObtenidasPorVuelo}{}\label{sec:MillasObtenidasPorVuelo}
\auxLargo{MillasObtenidasPorVuelo}
{Historial: \Tipo{\seq{Transacciones}}, idOrigen: \Tipo{\Z}}
{\Tipo{\R}}
{
	\sum\limits_{i=0}^{\norm{Historial} - 1}
	\IfThenElse{
		\hyperref[sec:EsVuelo]{EsVuelo}(Historial[i],
		idOrigen)
	}{
		Historial[i].Millas
	}{
		0
	}
}

\vspace{1cm}

\vspace{1cm}
\phantomsection\hypertarget{sec:MillasDescontadasPorCanjeo}{}\label{sec:MillasDescontadasPorCanjeo}
\auxLargo{MillasDescontadasPorCanjeo}
{Historial: \Tipo{\seq{Transacciones}}, idOrigen: \Tipo{\Z}}
{\Tipo{\R}}
{
	\sum\limits_{i=0}^{\norm{Historial} - 1}
	\IfThenElse{
		\hyperref[sec:EsCanjeo]{EsCanjeo}(Historial[i],
		idOrigen)
	}{
		Historial[i].Millas
	}{
		0
	}
}
\vspace{1cm}
\phantomsection\hypertarget{sec:EsVuelo}{}\label{sec:EsVuelo}
\predLargo{EsVuelo}
{Transaccion: \Tipo{Struct}, idOrigen: \Tipo{\Z}}
{
	Transaccion.UsuarioOrigen = idOrigen \land
	\\
	Transaccion.TipoDeTransaccion = 0 \land
	\\
	Transaccion.UsuarioDestino = -1
}

\vspace{1cm}


\phantomsection\hypertarget{sec:EsCanjeo}{}\label{sec:EsCanjeo}
\predLargo{EsCanjeo}
{Transaccion: \Tipo{struct}, idOrigen: \Tipo{\Z}}
{
	Transaccion.UsuarioOrigen = idOrigen \land
	\\
	Transaccion.TipoDeTransaccion = 2 \land
	\\
	Transaccion.UsuarioDestino = -1
}

\comentario{Es similar a EsTransferencia pero aca no importa el destino, unicamente chequeo que exista el usuario de origen.}
\\
\phantomsection\hypertarget{sec:EsTransferenciaEnviada}{}\label{sec:EsTransferenciaEnviada}
\predLargo{EsTransferenciaEnviada}
{Transaccion: \Tipo{Transacciones}, idOrigen: \Tipo{\Z}, Usuarios: \Tipo{\dict{\Z}{Categoria}}}
{
	Transaccion.UsuarioOrigen = idOrigen \land
	\\
	Transaccion.TipoDeTransaccion = 1 \land
	\\
	Transaccion.UsuarioDestino \in claves(Usuarios)
}

\vspace{1cm}
\phantomsection\hypertarget{sec:MillasPerdidas}{}\label{sec:MillasPerdidas}
\auxLargo{MillasPerdidas}
{Historial: \Tipo{\seq{Transacciones},idUsuario: \Tipo{\Z},Usuarios: \Diccionario{\Tipo{\Z}}{categoria}}}
{\Tipo{\R}}
{\sum\limits_{i=0}^{\norm{Historial}-1}
	\IfThenElse
	{\hyperref[sec:EsTransferenciaEnviada]{EsTransferenciaEnviada}(Historial[i],idUsuario,Usuarios) \lor
		\\
		\hyperref[sec:EsCanjeo]{EsCanjeo}(Historial[i],idUsuario)}
	{Historial[i].Millas}
	{0}
}

\phantomsection\hypertarget{sec:NoHayRepetidos}{}\label{sec:NoHayRepetidos}
\predLargo{NoHayRepetidos}{res: \Tipo{\seq{\tupla{\Z,\Z}}}}
    {
        \paraTodo{t}{\tupla{\Tipo{\Z},\Tipo{\Z}}}{t \in res \implicaLuego \hyperref[sec:cantApariciones]{cantApariciones}(t,res)=1}
    }

    \phantomsection\hypertarget{sec:cantApariciones}{}\label{sec:cantApariciones}
    \auxLargo{cantApariciones}{t:\Tipo{\tupla{\Tipo{\Z,\Z}}},res: \Tipo{\seq{\Tipo{\Z},\Tipo{\Z}}}}{\Tipo{\Z}}{
        \sum\limits_{0}^{\norm{res}} \IfThenElse{t \in res}{1}{0}
    }


    \comentario{Cual es la forma correcta de hacer esto? algo asi pero tengo que indicar el orden en el que estoy poniendo las cosas?}
    \\
    \phantomsection\hypertarget{sec:NoHayPermutacion}{}\label{sec:NoHayPermutacion}
    \predLargo{NoHayPermutacion}{res}{
        \paraTodo{t}{\Tipo{\tupla{\Tipo{\Z},\Tipo{\Z}}}}{t \in res \implicaLuego \tupla{{t_1},t_{0}} \notin res}
    }

    \vspace{0.5cm}
    \phantomsection\hypertarget{sec:TieneReciprocoHistorial}{}\label{sec:TieneReciprocoHistorial}
    \predLargo{TieneReciprocoHistorial}{Historial: \Tipo{\seq{Transacciones}},res: \seq{\tupla{\Tipo{\Z}, \Tipo{\Z}}}}
    {\paraTodo{x}{\tupla{\Tipo{\Z},\Tipo{\Z}}}{x \in res \implica \hyperref[sec:HayReciproco]{HayReciproco}(Historial,res)}}

    \vspace{0.5cm}
    \phantomsection\hypertarget{sec:HayReciproco}{}\label{sec:HayReciproco}
    \predLargo{HayReciproco}{Historial: \Tipo{\seq{Transacciones}},t: \Tipo{\Tipo{\tupla{\Tipo{\Z},\Tipo{\Z}}}}}
    {
        \existe{i}{\Tipo{\Z}}{0 \le i < \norm{Historial} \yLuego
        Historial[i].TipoDeTransaccion = 1 \yLuego 
        \\
        (Historial[i].usuarioOrigen = t_{1} \land Historial[i].usuarioDestino = t_{0})}
    }

    \vspace{0.5cm}
    \noindent \comentario{Revisar}
    \\
    \phantomsection\hypertarget{sec:SonParejaExclusiva}{}\label{sec:SonParejaExclusiva}
    \predLargo{SonParejaExclusiva}{Historial: \Tipo{\seq{Transacciones}},res: \seq{\tupla{\Tipo{\Z}, \Tipo{\Z}}}}
    {
        \paraTodo{t}{\tupla{\Tipo{\Z},\Tipo{\Z}}}{ t \in res \implicaLuego
        \neg (\hyperref[sec:HayOtroUsuario]{HayOtroUsuario}(t0,t1,Historial) \lor
        \\
        \hyperref[sec:HayOtroUsuario]{HayOtroUsuario}(t1,t0,Historial))
        }
    }

    \vspace{0.5cm}
    \phantomsection\hypertarget{sec:HayOtroUsuario}{}\label{sec:HayOtroUsuario}
    \predLargo{HayOtroUsuario}{t0,t1,Historial}
    {
        \paraTodo{i}{\Tipo{\Z}}{ 0 \le i < \norm{Historial} \yLuego Historial[i].UsuarioOrigen = t0 \implicaLuego 
        \\ 
        Historial[i].UsuarioDestino \neq t1}
    }


\end{tad}

\end{document}
